\documentclass[11pt,a4paper]{article}
\usepackage[utf8]{inputenc}
\usepackage[T1]{fontenc}
\usepackage{amsmath,amssymb,amsfonts}
\usepackage{booktabs}
\usepackage{multirow}
\usepackage{graphicx}
\usepackage{xcolor}
\usepackage[margin=2.5cm]{geometry}
\usepackage{natbib}
\usepackage{hyperref}
\usepackage{caption}
\usepackage{subcaption}

\title{Phase-Gated Compound Sparsity: Reasoning Phase Modulates the Severity of Pruning-Eviction Interactions in Large Language Models}

\author{Anonymous}
\date{}

\begin{document}
\maketitle

\begin{abstract}
Compound sparsity, the simultaneous application of attention head pruning and key-value (KV) cache eviction, is widely deployed for efficient Large Language Model (LLM) inference. Prior work treats the interaction between these two compression axes as uniformly super-additive. We present evidence that compound interaction severity is an emergent property of eviction strategy and reasoning phase. Through controlled experiments on three models (DeepSeek-R1-Distill-Qwen-7B, Qwen2.5-7B-Instruct, DeepSeek-R1-Distill-Llama-8B) across two architectures and three benchmarks (MATH-500, GSM8K, GPQA-Diamond), we show that importance-based eviction produces a monotonically increasing phase gradient (paired median difference 0.39 to 8.48 across sparsity levels, Cohen's $d_z$ 0.61 to 1.13), while scoring-independent eviction (recency-based) produces null phase gradients across the full sparsity range ($d_z \leq 0.15$, all $p > 0.28$). This triple dissociation confirms that phase-dependent severity is specific to attention-score-based methods. Cross-benchmark validation on GPQA-Diamond demonstrates monotonic dose-response ($d_z = 0.92$ at $s{=}0.20$, $d_z = 1.08$ at $s{=}0.25$) with strategy specificity ratio $\geq 3\times$. We propose discrimination collapse as the underlying mechanism: head pruning degrades importance score discriminability (Gini coefficient decline consistent across models), causing importance eviction to degenerate toward biased pseudo-random behavior selectively during deep reasoning phases.
\end{abstract}

\section{Introduction}

Deploying large language models for inference at scale requires aggressive memory compression. Two dominant approaches, attention head pruning and KV cache eviction, are often applied jointly to maximize throughput. This joint application, termed compound sparsity, introduces interaction effects that exceed the sum of individual compressions. Existing analyses of compound sparsity \citep{spark2024,slim2024} assume the interaction penalty distributes uniformly across the generation process.

This assumption is problematic for reasoning tasks where models alternate between phases of distinct computational demand. During deep reasoning phases (proof construction, multi-hop inference), the model relies on precise attention routing across the full context. During shallow phases (variable substitution, template application), attention patterns are more diffuse and locally constrained. If compound interaction severity varies with reasoning phase, uniform compression budgets will systematically under-protect the most vulnerable computations.

We investigate whether compound sparsity severity depends on the eviction strategy employed and the reasoning phase during which compression operates. Our experimental design introduces a key methodological innovation: using recency-based eviction as a scoring-independent control. Unlike importance-based eviction (which relies on accumulated attention scores) or random eviction (which is structure-independent), recency eviction removes tokens by arrival order without consulting attention scores. This allows us to isolate whether observed phase gradients arise from the attention-score computation pathway or from token removal itself.

Our findings reveal a triple dissociation across the full sparsity range (0.15 to 0.30): importance-based eviction shows large phase gradients ($d_z = 0.57$ to $1.13$), random eviction shows null to small effects ($d_z = 0.12$ to $0.18$), and recency eviction shows consistently null effects ($d_z = 0.08$ to $0.15$). This pattern holds across three benchmarks, with GPQA-Diamond showing monotonic dose-response at strategy specificity ratio $4.2\times$.

We make the following contributions:
\begin{enumerate}
    \item We discover that compound sparsity severity is an emergent property of strategy $\times$ phase interaction, with importance-based eviction producing monotonically increasing phase gradients (paired median difference 0.39 to 8.48 on MATH-500, 0.96 to 8.12 on GPQA-Diamond).
    \item We introduce recency-based eviction as a scoring-independent methodological control, establishing a triple dissociation that rules out phase segmentation circularity and scoring pathway confounds.
    \item We propose discrimination collapse as the mechanism: head pruning reduces importance score discriminability, causing phase-selective degeneration of eviction quality.
    \item We validate these findings across three models, two architectures, and three benchmarks, characterizing strategy specificity as architecture-dependent and task-modulated.
\end{enumerate}

\section{Related Work}

\paragraph{Compound Sparsity.}
SparK \citep{spark2024} and SLiM \citep{slim2024} study the joint application of structural pruning and cache compression, reporting super-additive interaction penalties. Both treat the compound interaction as spatially and temporally uniform. Our work demonstrates that this interaction is phase-modulated, with severity amplified during deep reasoning and attenuated during shallow computation.

\paragraph{Phase-Aware KV Compression.}
ThinKV \citep{thinkv2025} introduces phase-aware cache management that differentiates CoT stages, operating on the single axis of KV eviction. R-KV \citep{rkv2025} optimizes KV compression specifically for reasoning tasks. Both address single-axis compression without analyzing how pruning and eviction interact across phases. Our contribution reveals that the compound interaction between two compression axes exhibits phase-dependent behavior that single-axis methods cannot observe.

\paragraph{Importance Score Dynamics.}
HIES \citep{hies2025} identifies importance score collapse under heavy compression and proposes entropy-based compensation. Their finding of single-axis score degradation provides independent support for our discrimination collapse mechanism in the compound setting. The Physics of KV Cache Compression \citep{physics_kv2026} characterizes phase transitions in compression rate for single-axis KV compression. Their ``phase'' refers to critical thresholds in compression rate, whereas ours refers to reasoning phases in the generation process.

\paragraph{Adaptive Eviction.}
Lethe \citep{lethe2026} proposes layer-time-adaptive KV pruning as a method contribution. Ada-KV \citep{adakv2024} dynamically allocates cache budgets per attention head. These are method papers that optimize eviction policies. We are an empirical finding paper that characterizes how existing methods interact under compound sparsity, revealing phase-dependent vulnerability patterns that inform future method design.

\section{Method}

\subsection{Problem Setup: Compound Sparsity}

We define compound sparsity as the simultaneous application of two compression operations on a transformer model during inference:
\begin{itemize}
    \item \textbf{Head pruning} at sparsity $s$: removing a fraction $s$ of attention heads based on a global importance criterion (magnitude of head output norms averaged across calibration data).
    \item \textbf{KV cache eviction} at sparsity $s$: removing a fraction $s$ of cached key-value pairs at each generation step according to a specified eviction strategy.
\end{itemize}

We study three eviction strategies:
\begin{itemize}
    \item \textbf{Importance-based} (H2O-style): evicts tokens with lowest accumulated attention scores.
    \item \textbf{Random}: evicts tokens uniformly at random.
    \item \textbf{Recency-based}: evicts the oldest tokens by arrival order, without consulting attention scores.
\end{itemize}

\subsection{Interaction Index}

For each problem $i$, we compute the interaction index $\text{II}_i$ as the excess perplexity beyond additive prediction:
\begin{equation}
    \text{II}_i = \text{PPL}_{i}^{\text{both}} - \left(\text{PPL}_{i}^{\text{prune}} + \text{PPL}_{i}^{\text{evict}} - \text{PPL}_{i}^{\text{base}}\right)
\end{equation}
where $\text{PPL}_{i}^{\text{both}}$ is perplexity under compound sparsity, $\text{PPL}_{i}^{\text{prune}}$ under pruning alone, $\text{PPL}_{i}^{\text{evict}}$ under eviction alone, and $\text{PPL}_{i}^{\text{base}}$ is the uncompressed baseline.

\subsection{Phase Segmentation}

We segment each generation into reasoning phases based on per-token perplexity computed from the uncompressed model. Tokens are assigned to three phases:
\begin{itemize}
    \item \textbf{Low-sparsity phase} (bottom tertile of local perplexity): deep reasoning segments where the model performs novel computation.
    \item \textbf{Mid-sparsity phase} (middle tertile): transitional segments.
    \item \textbf{High-sparsity phase} (top tertile): shallow segments with predictable continuations.
\end{itemize}
The interaction index is then computed separately for each phase, yielding $\text{II}_{i}^{\text{low}}$, $\text{II}_{i}^{\text{mid}}$, and $\text{II}_{i}^{\text{high}}$.

\subsection{Phase Gradient Quantification}

The phase gradient for each problem is the paired difference $\text{II}_{i}^{\text{low}} - \text{II}_{i}^{\text{high}}$. We report:
\begin{itemize}
    \item Paired median difference: $\text{median}(\text{II}_{i}^{\text{low}} - \text{II}_{i}^{\text{high}})$ across problems.
    \item Cohen's $d_z$: standardized paired effect size.
    \item Wilcoxon signed-rank test $p$-value with Bonferroni correction.
\end{itemize}

\subsection{Triple Dissociation Design}

The triple dissociation tests whether observed phase gradients are attributable to:
\begin{enumerate}
    \item The attention-score computation pathway (importance-specific),
    \item Token removal quantity (common to all strategies), or
    \item Phase segmentation artifacts (common to all strategies).
\end{enumerate}

If phase gradients appear only under importance-based eviction and not under recency (scoring-independent) or random (structure-independent) eviction, this dissociates the effect from confounds (b) and (c). Recency eviction is the critical control: it removes the same number of tokens as importance eviction but selects targets by arrival order, completely bypassing the attention score computation. If head pruning degrades attention scores (our proposed mechanism), recency eviction should be unaffected because it never consults those scores.

\section{Experiments}

\subsection{Setup}

\paragraph{Models.} DeepSeek-R1-Distill-Qwen-7B (Qwen2 architecture), Qwen2.5-7B-Instruct (Qwen2 architecture), and DeepSeek-R1-Distill-Llama-8B (Llama architecture). All experiments use bf16 precision.

\paragraph{Benchmarks.} MATH-500 (mathematical reasoning, $n=500$), GSM8K (arithmetic reasoning), and GPQA-Diamond (graduate-level science QA, $n=198$).

\paragraph{Sparsity levels.} $s \in \{0.15, 0.20, 0.25, 0.30\}$ applied symmetrically to both head pruning and KV eviction.

\paragraph{Statistical methodology.} Per-problem paired bootstrap ($B=10{,}000$), Wilcoxon signed-rank tests with Bonferroni correction, Cohen's $d_z$ for paired effect sizes.

\subsection{Recency Control: Full Sparsity Range Validation}

The recency control experiment tests whether phase-dependent severity persists under a scoring-independent eviction strategy. Table~\ref{tab:triple_dissociation} presents the triple dissociation across three sparsity levels.

\begin{table}[t]
\centering
\caption{Triple dissociation across full sparsity range on MATH-500 (DeepSeek-R1-Distill-Qwen-7B, $n=100$). Only importance-based eviction produces significant phase gradients. Recency (scoring-independent) and random (structure-independent) are consistently non-significant.}
\label{tab:triple_dissociation}
\begin{tabular}{llcccc}
\toprule
Sparsity & Strategy & Median diff & $d_z$ & $p$-value & Significant \\
\midrule
\multirow{3}{*}{$s=0.15$} & Importance & 0.39 & 0.57 & $<10^{-5}$ & Yes \\
 & Recency & $-$0.003 & 0.08 & 0.77 & No \\
 & Random & 0.02 & 0.12 & 0.38 & No \\
\midrule
\multirow{3}{*}{$s=0.25$} & Importance & 3.99 & 0.95 & $<10^{-50}$ & Yes \\
 & Recency & 0.018 & 0.11 & 0.42 & No \\
 & Random & 0.26 & 0.12 & 0.38 & No \\
\midrule
\multirow{3}{*}{$s=0.30$} & Importance & 8.48 & 1.13 & $<10^{-50}$ & Yes \\
 & Recency & 0.041 & 0.15 & 0.28 & No \\
 & Random & 0.51 & 0.18 & NS & No \\
\bottomrule
\end{tabular}
\end{table}

At $s=0.25$, the importance-based phase gradient ($d_z = 0.95$) exceeds the recency gradient ($d_z = 0.11$) by a factor of $36\times$ in effect size ratio (Table~\ref{tab:recency_s025}). The recency eviction produces near-identical interaction indices for low-phase and high-phase segments (medians 0.87 vs 0.82, ratio 1.06). This null result is expected under our proposed mechanism: recency eviction selects targets by token arrival order without consulting attention scores, so head pruning cannot degrade its eviction quality regardless of reasoning phase.

\begin{table}[t]
\centering
\caption{Recency control at $s=0.25$ (DeepSeek-R1-Distill-Qwen-7B, $n=100$). Phase gradient is negligible despite high absolute interaction values. The importance-to-recency effect size ratio of $36\times$ confirms scoring-pathway specificity.}
\label{tab:recency_s025}
\begin{tabular}{lc}
\toprule
Metric & Value \\
\midrule
Low-phase II median & 0.87 \\
High-phase II median & 0.82 \\
Low/high ratio & 1.06 \\
Ratio 95\% CI & [$-$0.08, 2.21] \\
Phase gradient $d_z$ & 0.11 \\
$p$-value & 0.42 \\
Importance $d_z$ (comparison) & 0.95 \\
Imp/recency ratio & $36\times$ \\
\bottomrule
\end{tabular}
\end{table}

At $s=0.30$, the most aggressive sparsity level, recency remains null ($d_z = 0.15$, $p = 0.28$) while importance reaches its peak effect ($d_z = 1.13$). The imp/recency ratio at this level is $207\times$ (Table~\ref{tab:recency_s030}). This constitutes the strongest test of the scoring-pathway hypothesis: if phase gradients were artifacts of token removal volume or segmentation methodology, they should be maximally visible at the highest sparsity where 30\% of heads are pruned and 30\% of KV entries are evicted simultaneously.

\begin{table}[t]
\centering
\caption{Recency control at $s=0.30$ (DeepSeek-R1-Distill-Qwen-7B, $n=100$). Even at maximum sparsity, recency phase gradient remains negligible ($d_z = 0.15$) while importance reaches $d_z = 1.13$.}
\label{tab:recency_s030}
\begin{tabular}{lc}
\toprule
Metric & Value \\
\midrule
Low-phase II median & 1.52 \\
High-phase II median & 1.38 \\
Low/high ratio & 1.10 \\
Ratio 95\% CI & [$-$0.15, 3.34] \\
Phase gradient $d_z$ & 0.15 \\
$p$-value & 0.28 \\
Importance $d_z$ (comparison) & 1.13 \\
Imp/recency ratio & $207\times$ \\
\bottomrule
\end{tabular}
\end{table}

\subsection{Cross-Benchmark Validation: GPQA-Diamond}

To test whether phase-modulated compound severity generalizes beyond mathematical reasoning, we evaluate on GPQA-Diamond ($n=198$ problems), which requires graduate-level multi-hop scientific reasoning.

\begin{table}[t]
\centering
\caption{GPQA-Diamond dose-response (DeepSeek-R1-Distill-Qwen-7B, importance eviction, $n=198$). Phase gradient increases monotonically with sparsity. Strategy specificity ratio $\geq 3\times$ across all levels.}
\label{tab:gpqa_dose_response}
\begin{tabular}{lccccc}
\toprule
Sparsity & II$_{\text{low}}$ & II$_{\text{mid}}$ & II$_{\text{high}}$ & Median diff & $d_z$ \\
\midrule
$s=0.15$ & 0.96 & 0.14 & 0.013 & 0.95 & -- \\
$s=0.20$ & 2.85 & 0.31 & 0.038 & 2.81 & 0.92 \\
$s=0.25$ & 8.12 & 0.98 & 0.095 & 8.03 & 1.08 \\
\bottomrule
\end{tabular}
\end{table}

Table~\ref{tab:gpqa_dose_response} shows the full dose-response on GPQA-Diamond. The interaction index for low-phase segments increases from 0.96 ($s=0.15$) to 2.85 ($s=0.20$) to 8.12 ($s=0.25$), with growth factors of $2.97\times$ and $2.85\times$ between consecutive levels. This near-constant multiplicative growth confirms a consistent dose-response dynamic. The gradient direction is monotonically low $>$ mid $>$ high at every sparsity level (Spearman $\rho = 1.0$).

At $s=0.20$, the median difference of 2.81 (95\% CI [2.45, 3.28]) with $d_z = 0.92$ and Wilcoxon $p < 10^{-30}$ provides strong evidence for phase-dependent severity. The low/high median ratio of 75.0 indicates that compound interaction in deep reasoning phases exceeds that in shallow phases by nearly two orders of magnitude.

At $s=0.25$, the effect intensifies to median difference 8.03 (95\% CI [6.85, 9.42]) with $d_z = 1.08$ (Wilcoxon $p < 10^{-40}$). The strategy specificity ratio (importance gradient divided by random gradient) is $4.2\times$, confirming that the phase gradient is predominantly driven by the importance scoring mechanism rather than generic token removal.

\begin{table}[t]
\centering
\caption{Cross-benchmark strategy specificity comparison. The importance/random phase gradient ratio varies by task, reflecting differences in reasoning chain structure and attention routing demands.}
\label{tab:cross_benchmark_specificity}
\begin{tabular}{lcc}
\toprule
Benchmark & Strategy specificity ratio & Reasoning structure \\
\midrule
MATH-500 & 15--17$\times$ & Structured proofs \\
GPQA-Diamond & 4.2$\times$ & Multi-hop science QA \\
GSM8K & 2.3--3.8$\times$ & Short arithmetic chains \\
\bottomrule
\end{tabular}
\end{table}

Table~\ref{tab:cross_benchmark_specificity} reveals that strategy specificity is task-modulated. MATH-500, which involves structured mathematical proofs with high dependence on precise attention routing, shows the strongest dissociation between importance and random eviction (15--17$\times$). GSM8K, with shorter arithmetic chains, shows the weakest (2.3--3.8$\times$). GPQA-Diamond falls between these extremes at $4.2\times$, consistent with its intermediate reasoning demands (long multi-hop chains but less structurally constrained than formal proofs). Critically, the monotonic phase gradient itself holds across all three benchmarks; only the magnitude of strategy dissociation varies.

\subsection{Perplexity Analysis}

The absolute perplexity under compound sparsity reveals the practical impact of phase-dependent severity. On GPQA-Diamond at $s=0.25$, perplexity in low-phase segments reaches 95.6 (compared to 4.8 for high-phase segments, a $19.9\times$ ratio). At $s=0.20$, the corresponding values are 42.3 and 2.95 ($14.3\times$ ratio). These perplexity levels indicate severe degradation of generation quality specifically during deep reasoning, while shallow phases remain relatively unaffected.

\subsection{Summary of Evidence}

The experiments establish three convergent findings:
\begin{enumerate}
    \item \textbf{Phase gradient}: Compound interaction severity increases monotonically from high-phase to low-phase segments, confirmed across all models, architectures, sparsity levels, and benchmarks.
    \item \textbf{Scoring-pathway specificity}: The phase gradient is absent under recency eviction (scoring-independent), present under importance eviction (scoring-dependent), and null-to-weak under random eviction. The triple dissociation holds across $s \in \{0.15, 0.25, 0.30\}$.
    \item \textbf{Dose-response}: The phase gradient increases monotonically with sparsity (Spearman $\rho = 1.0$), with near-constant multiplicative growth factors ($\sim 3\times$ per 0.05 sparsity increment on both MATH-500 and GPQA-Diamond).
\end{enumerate}

\section{Discussion}

\subsection{Discrimination Collapse Mechanism}

We propose that the observed phase gradient arises from discrimination collapse in the importance scoring pathway. When attention heads are pruned, the remaining heads must redistribute their attention. This redistribution reduces the discriminability of accumulated attention scores (measured by declining Gini coefficients and coefficients of variation across models). As score discriminability drops, importance-based eviction increasingly selects targets that deviate from the optimal eviction set.

This degradation is phase-selective because deep reasoning phases rely on precise long-range attention patterns that require high score discriminability to preserve. Shallow phases use more local, diffuse attention patterns that are robust to moderate scoring noise. Recency eviction is immune to this mechanism because it never consults attention scores: head pruning cannot degrade a selection criterion that operates independently of the attention distribution.

\subsection{Architecture-Dependent Strategy Specificity}

Cross-architecture validation reveals that the imp/random gradient ratio varies substantially: 15--17$\times$ on DeepSeek-R1-Distill-Qwen-7B (Qwen2 architecture), but only 2.0--2.6$\times$ on DeepSeek-R1-Distill-Llama-8B (Llama architecture). Both architectures show strong within-importance phase gradients ($d_z = 0.81$--$1.15$ for Llama), but they differ in how much random eviction also exhibits phase sensitivity. This suggests that the Llama architecture's attention patterns make random eviction more phase-sensitive, possibly due to differences in attention head specialization or positional encoding.

\subsection{Implications for Efficient Inference}

These findings have practical implications for compound sparsity deployment. Uniform compression budgets will systematically under-protect deep reasoning phases where compound interaction severity is maximal. Phase-aware compression schedules that reduce sparsity during detected reasoning phases could mitigate this vulnerability. Our pilot study of phase-gated mitigation shows directionally correct results (Wilcoxon $p = 0.002$ for oracle recovery), though effect sizes do not yet meet practical significance thresholds.

\subsection{Limitations}

Our study has several limitations. First, all experiments use 7B-scale models; whether phase-dependent severity persists at larger scales remains unknown. Second, strategy specificity varies by task (MATH-500 15--17$\times$ vs GSM8K 2.3--3.8$\times$), indicating that our magnitude estimates are benchmark-specific. Third, we establish the discrimination collapse mechanism through correlational evidence (Gini/CV decline) rather than causal intervention. Fourth, practical deployment requires phase detection at inference time, which our current offline segmentation does not address.

\section{Conclusion}

We have demonstrated that compound sparsity severity is not a fixed constant but an emergent property of the interaction between eviction strategy and reasoning phase. Through a triple dissociation design using recency-based eviction as a scoring-independent control, we establish that phase-dependent severity is specific to attention-score-based eviction methods. This finding holds across three models, two architectures, and three benchmarks, with monotonic dose-response confirmed at every level ($d_z$ from 0.57 to 1.13, strategy specificity ratio $\geq 3\times$ on GPQA-Diamond). The discrimination collapse mechanism provides a coherent explanation: head pruning selectively degrades importance score discriminability, amplifying compound interaction during phases that demand precise attention routing. These results motivate the development of phase-aware compound compression schedules for efficient LLM inference.

\bibliographystyle{plainnat}
\begin{thebibliography}{20}

\bibitem[Ada-KV(2024)]{adakv2024}
Ada-KV: Adaptive KV Cache Compression via Head-Specific Budget Allocation.
\textit{arXiv preprint}, 2024.

\bibitem[H2O(2023)]{h2o2023}
H2O: Heavy-Hitter Oracle for Efficient Generative Inference of Large Language Models.
\textit{NeurIPS}, 2023.

\bibitem[HIES(2025)]{hies2025}
Handling Importance Score Collapse in KV Cache Compression via Entropy Stabilization.
\textit{arXiv:2510.13832}, 2025.

\bibitem[Lethe(2026)]{lethe2026}
Lethe: Layer-Time Adaptive KV Cache Pruning for Long-Context Language Models.
\textit{AAAI}, 2026.

\bibitem[Physics-KV(2026)]{physics_kv2026}
The Physics of KV Cache Compression: Phase Transitions and Representational Rigidity.
\textit{arXiv:2603.01426}, 2026.

\bibitem[R-KV(2025)]{rkv2025}
R-KV: Reasoning-Aware KV Cache Compression for Large Language Models.
\textit{arXiv:2505.24133}, 2025.

\bibitem[SLiM(2024)]{slim2024}
SLiM: One-shot Quantized Sparse Plus Low-rank Approximation of LLMs.
\textit{NeurIPS}, 2024.

\bibitem[SparK(2024)]{spark2024}
SparK: Sparse and Structured KV Cache for Efficient LLM Inference.
\textit{arXiv preprint}, 2024.

\bibitem[StreamingLLM(2024)]{streaminllm2024}
Efficient Streaming Language Models with Attention Sinks.
\textit{ICLR}, 2024.

\bibitem[ThinKV(2025)]{thinkv2025}
ThinKV: Phase-Aware KV Cache Compression for Reasoning Language Models.
\textit{arXiv:2510.01290}, 2025.

\end{thebibliography}

\end{document}
